\documentclass[11pt]{article}
\input{style-sujet}

\begin{document}

\entetesujet{Sujet bac 2010 -- Série D}

% ====================== EXERCICE 1 ======================
\exo{1}{4 points}

Une urne contient deux boules blanches et trois boules noires toutes indiscernables au toucher. On tire simultanément et au hasard deux boules de l'urne et on note leur couleur.

On définit sur l'univers $\Omega$ de cette expérience aléatoire, la variable aléatoire réelle $X$ par :
\begin{itemize}
  \item $X = -1$, si les deux boules tirées sont blanches ;
  \item $X = 0$, si l'une est blanche et l'autre est noire ;
  \item $X = 1$, si les deux boules tirées sont noires.
\end{itemize}

\begin{enumerate}
  \item Déterminer la loi de probabilité de $X$.
  \item Calculer l'espérance mathématique de $X$.
  \item \begin{enumerate}[label=\alph*.]
    \item Définir la fonction de répartition $F$ de $X$.
    \item Tracer la courbe représentative de $F$ dans un repère orthogonal (on prendra 1 cm en abscisse et 5 cm en ordonnée pour unités graphiques).
  \end{enumerate}
\end{enumerate}

% ====================== EXERCICE 2 ======================
\exo{2}{4 points}

Dans le plan complexe muni d'un repère orthonormal direct $(O,\vec{u},\vec{v})$, on considère les nombres complexes $z_1 = \sqrt{3} + i$ ; $z_2 = -\sqrt{3} + i$ et $z_3 = -2i$.

\begin{enumerate}
  \item Écrire une équation de degré 3 dont $z_1$, $z_2$ et $z_3$ sont solutions.
  \item Écrire les nombres complexes $z_1$, $z_2$ et $z_3$ sous la forme trigonométrique.
  \item \begin{enumerate}[label=\alph*.]
    \item Placer les points $A$, $B$ et $C$ d'affixes respectives $z_1$, $z_2$ et $z_3$ dans le repère $(O,\vec{u},\vec{v})$.
    \item Montrer qu'une mesure de chacun des angles du triangle $ABC$ est $\dfrac{\pi}{3}$ radians.
    \item En déduire la nature du triangle $ABC$.
  \end{enumerate}
\end{enumerate}

% ====================== PROBLÈME ======================
\prob{12 points}

\partie{A}
Soit $g$ la fonction de la variable réelle $x$ définie sur $]-\infty\,;0[$ par :
\[ g(x) = 2\ln(-x) \]

\begin{enumerate}
  \item Étudier les variations de $g$, puis dresser son tableau de variation.
  \item Calculer $g(-1)$ et en déduire le signe de $g(x)$ sur $]-\infty\,;0[$.
\end{enumerate}

\partie{B}
On considère la fonction numérique $f$ de la variable réelle $x$ définie par :
\[ f(x) = \begin{cases} -2x + 1 + 2x\ln|x| & \text{si } x < 0 \\[0.5ex] (x+2)e^{-x} - 1 & \text{si } x \ge 0 \end{cases} \]
On désigne par $(C)$, la courbe représentative de $f$ dans un repère orthonormé $(O,\vi,\vj)$ d'unité graphique : 1 cm.

\begin{enumerate}
  \item Déterminer l'ensemble de définition de la fonction $f$.
  \item \begin{enumerate}[label=\alph*.]
    \item Étudier la continuité et la dérivabilité de $f$ au point $x = 0$.
    \item Pour $x \in ]-\infty\,;0[$, exprimer $f'(x)$ en fonction de $g(x)$.
  \end{enumerate}
  \item Étudier les variations de $f$ et dresser son tableau de variation.
  \item Montrer qu'il existe deux solutions et deux seulement $\alpha$ et $\beta$ de l'équation $f(x) = 0$ vérifiant les inégalités suivantes :
  \[ 1 < \alpha < \dfrac{3}{2} \quad \text{et} \quad -4 < \beta < -3 \]
  (on ne cherchera pas à calculer $\alpha$ et $\beta$).
  \item Étudier les branches infinies à $(C)$.
  \item Tracer la courbe $(C)$.
  \item $\alpha$ désigne le réel tel que : $1 < \alpha < \dfrac{3}{2}$.
  \begin{enumerate}[label=\alph*.]
    \item Calculer l'aire $\mathcal{A}(\alpha)$ de l'ensemble des points $M$ de coordonnées $(x, y)$ tels que :
    \[ \begin{cases} 0 \le x \le \alpha \\ 0 \le y \le f(x) \end{cases} \]
    \item Calculer la limite de $\mathcal{A}(\alpha)$ lorsque $\alpha$ tend vers 0.
  \end{enumerate}
\end{enumerate}

\partie{C}
Soit $h$ la restriction de $f$ à l'intervalle $I = ]-\infty\,;-1]$.

\begin{enumerate}
  \item Montrer que $h$ définit une bijection de $I$ vers un intervalle $J$ à déterminer. On note $h^{-1}$ la réciproque de $h$.
  \item Dresser le tableau de variation de $h^{-1}$.
  \item Tracer $(H)$ la courbe représentative de $h^{-1}$ dans le même repère que $(C)$.
\end{enumerate}

\end{document}
